% =====================================================================
%   Подведение итогов главы
% =====================================================================
\section*{Подведение итогов главы}
\addcontentsline{toc}{section}{Подведение итогов главы}
\markboth{Глава \thechapter. Элементы теории чисел и шифрование}{Подведение итогов главы}

\begin{tcolorbox}[
  enhanced, sharp corners, boxrule=0pt,
  colback=prosvCream, leftrule=4pt, colframe=prosvBlue,
  left=12pt, right=12pt, top=10pt, bottom=10pt,
  before skip=0.5em, after skip=1.5em,
  title={\textbf{\color{prosvBlue}\large Главное в этой главе}},
  fonttitle=\bfseries
]
\textbf{Теория чисел.} Простые числа, остатки от деления, сравнения по модулю --- основы модулярной арифметики. Алгоритм Евклида и его расширенная версия позволяют не только находить НОД, но и решать модулярные линейные уравнения. Быстрое возведение в степень за $O(\log k)$ операций позволяет работать с гигантскими числами; та же идея даёт логарифмический алгоритм для чисел Фибоначчи через возведение матрицы $2{\times}2$ в степень. Малая теорема Ферма $a^{p-1} \equiv 1 \pmod p$ --- ключевое утверждение, обеспечивающее корректность RSA. Простых чисел много (закон $\pi(N) \sim N/\ln N$), и проверка простоты быстра (тест Миллера--Рабина, AKS), но \emph{факторизация} составного числа на множители --- задача экспоненциальной сложности.

\textbf{История криптографии.} От шифра Цезаря --- к симметричным шифрам с одним общим ключом --- к проблеме доставки ключа. Идея \term{односторонней функции} (легко в одну сторону, очень сложно в обратную) и \term{функции с секретом} (трапдор) даёт возможность строить асимметричные шифры с открытым и закрытым ключами.

\textbf{RSA и Диффи--Хеллман.} RSA опирается на сложность факторизации $n = pq$: открытый ключ $(n, e)$ позволяет шифровать, закрытый $d = e^{-1} \bmod \Eul(n)$ --- расшифровывать. Корректность доказывается малой теоремой Ферма. Диффи--Хеллман --- протокол выработки общего ключа по открытому каналу: его стойкость основана на сложности задачи дискретного логарифма. Обе схемы уязвимы к атаке «человек посередине» без аутентификации.

\textbf{Применения RSA.} Та же конструкция, применённая в разных порядках, решает четыре задачи: шифрование сообщения, аутентификация пользователя через challenge-response, электронная цифровая подпись (подпись закрытым ключом, проверка открытым) и слепая подпись (основа протоколов электронного голосования).

\textbf{Хеширование.} Универсальное семейство Картера--Вегмана $h_a(x) = (a_1 x_1 + \dots + a_k x_k) \bmod p$ ($p$ --- простое) гарантирует, что вероятность коллизии любых двух ключей при случайном выборе $a$ не превосходит $1/p$. Ключевую роль играет существование и единственность обратного элемента по простому модулю --- результат расширенного алгоритма Евклида.

\textbf{Вероятностные счётчики.} HyperLogLog оценивает число различных элементов в потоке с точностью $\sim 1\%$, используя всего килобайты памяти. В сочетании с теорией графов это позволяет современным сетям подтверждать гипотезу Стенли Мильграма о малом мире: средняя длина «социальной» цепочки в крупных сетях --- около \emph{четырёх рукопожатий}.
\end{tcolorbox}

\subsection*{Что почитать дальше}
\addcontentsline{toc}{subsection}{Что почитать дальше}

\begin{itemize}
  \item \emph{Музыкантский А.\,И., Фурин В.\,В.} Лекции по криптографии. М.: МЦНМО, разные годы. --- Доступное изложение криптографических протоколов, в том числе электронного голосования; рекомендуется к разделу~1.3 и главе~2.
  \item \emph{Литвак Н.\,В., Райгородский А.\,М.} Кому нужна математика? Содержательное введение в математику и computer science. --- Главы 6, 7 и приложения к ним; материал о малых мирах, графах и вероятностных алгоритмах.
  \item \emph{Дасгупта С., Пападимитриу Х., Вазирани У.} Алгоритмы. --- Глава о хешировании, главы об алгоритмах теории чисел.
  \item \emph{Кнут Д.} Искусство программирования, том~2 (Получисленные алгоритмы). --- Классический фундаментальный труд по алгоритмам теории чисел.
  \item \emph{Виноградов И.\,М.} Основы теории чисел. --- Классический российский учебник теории чисел, доступный школьникам старших классов.
\end{itemize}

\subsection*{Большой итоговый проект}
\addcontentsline{toc}{subsection}{Большой итоговый проект}

В качестве длинного исследовательского проекта на каникулы или на год предлагаем выбрать одно из следующих направлений:

\begin{itemize}
  \item \textbf{«Своя» реализация RSA.} На любом удобном языке программирования (Python, Java, C++) реализуйте полный цикл: генерацию пары ключей $(e, n)$ и $d$ (с честным алгоритмом Миллера--Рабина для подбора простых), шифрование, расшифрование, цифровую подпись. Протестируйте на примерах из этой главы.
  \item \textbf{Атака на «слабый» RSA.} Возьмите малое $n$ (50--100 битов) и попытайтесь его факторизовать. Сравните скорость наивного перебора и метода Полларда $\rho$. Подумайте, при каких размерах $n$ ваш ноутбук «сдаётся».
  \item \textbf{Симуляция «малого мира».} Сгенерируйте случайный граф из 1000--10000 вершин по модели Уоттса--Строгаца. Вычислите средние длины кратчайших путей. Сравните с предсказанием теории.
  \item \textbf{Свой HyperLogLog.} Реализуйте упрощённый вариант HyperLogLog ($b = 8$ или $10$). Сгенерируйте поток из миллиона случайных идентификаторов с заранее известным числом уникальных. Оцените точность алгоритма.
  \item \textbf{Стойкость хеш-функции.} Возьмите криптографическую хеш-функцию (например, SHA-1 уже взломан) и реализуйте поиск коллизий простым перебором. Сравните с теоретической оценкой «парадокса дней рождения».
\end{itemize}
